\subsection{Дифференциал. Геометрический смысл производной и дифференциала.}

\subsubsection*{Дифференцируемость и дифференциал}

\begin{tcolorbox}[title=Определение]
    Пусть функция $f(x)$ определена в окрестности $U_{x_0}$ и $\exists A \in \mathbb{R}$.
    Функция $f(x)$ называется \textbf{дифференцируемой} в точке $x_0$, если её приращение можно представить в виде:
    \[ f(x) = f(x_0) + A(x-x_0) + \bar{o}(x-x_0) \quad \text{при } x \to x_0 \]
    
    Слагаемое $A(x-x_0)$ (или $A\Delta x$) называется \textbf{дифференциалом} функции $f(x)$ в точке $x_0$ и обозначается $df$:
    \[ df = A(x-x_0) = A\Delta x \]
\end{tcolorbox}

\textbf{Замечание.} Если $f(x) = x$, то $df = dx = 1 \cdot (x-x_0) = \Delta x$.
Поэтому дифференциал часто записывают как $df = A dx$ (где $x$ — независимая переменная).

\begin{tcolorbox}[title=Теорема (Связь производной и дифференциала)]
    Функция $f(x)$ имеет производную в точке $x_0$ $\iff$ $f(x)$ дифференцируема в точке $x_0$.
    При этом $A = f'(x_0)$.
\end{tcolorbox}

\begin{tcolorbox}[title=Доказательство]
    \textbf{1. ($\Rightarrow$)}
    Пусть существует $f'(x_0) = \lim_{x\to x_0} \frac{f(x)-f(x_0)}{x-x_0}$.
    По определению предела функции:
    \[ \frac{f(x)-f(x_0)}{x-x_0} = f'(x_0) + \alpha(x), \quad \text{где } \alpha(x) \to 0 \text{ при } x \to x_0 \]
    Умножим на $(x-x_0)$:
    \[ f(x) - f(x_0) = f'(x_0)(x-x_0) + \alpha(x)(x-x_0) \]
    Заметим, что $\alpha(x)(x-x_0) = o(x-x_0)$.
    \[ f(x) = f(x_0) + f'(x_0)(x-x_0) + o(x-x_0) \]
    Это и означает дифференцируемость с $A = f'(x_0)$.

    \textbf{2. ($\Leftarrow$)}
    Пусть $f(x)$ дифференцируема: $f(x) - f(x_0) = A(x-x_0) + o(x-x_0)$.
    Разделим на $(x-x_0) \neq 0$:
    \[ \frac{f(x)-f(x_0)}{x-x_0} = A + \frac{o(x-x_0)}{x-x_0} \]
    Перейдем к пределу при $x \to x_0$. Так как $\frac{o(x-x_0)}{x-x_0} \to 0$, получаем:
    \[ \lim_{x\to x_0} \frac{f(x)-f(x_0)}{x-x_0} = A \]
    Предел существует, значит производная существует и $f'(x_0) = A$.
\end{tcolorbox}

\textbf{Формулы:}
\begin{itemize}
    \item $df = f'(x_0)dx$
    \item $f(x_0) \approx \frac{df}{dx}$ (для малых приращений)
\end{itemize}

\textbf{Теорема 2.} Если $f(x)$ дифференцируема в точке $x_0 \Rightarrow f(x)$ непрерывна в точке $x_0$.
(Доказательство следует из формулы приращения: при $x \to x_0$ правая часть стремится к $f(x_0)$).

\begin{tcolorbox}[title=Доказательство]
    \begin{enumerate}
        \item По условию $f(x)$ дифференцируема в точке $x_0$. По определению это значит, что приращение функции $\Delta f = f(x_0 + \Delta x) - f(x_0)$ представимо в виде:
        \[ \Delta f = A \cdot \Delta x + o(\Delta x), \quad \text{при } \Delta x \to 0 \]
        (где $A = f'(x_0)$ — конечное число).

        \item Найдем предел приращения функции при $\Delta x \to 0$:
        \[ \lim_{\Delta x \to 0} \Delta f = \lim_{\Delta x \to 0} \left( A \cdot \Delta x + o(\Delta x) \right) \]

        \item Используя свойства пределов:
        \[ \lim_{\Delta x \to 0} (A \cdot \Delta x) = A \cdot 0 = 0 \]
        \[ \lim_{\Delta x \to 0} o(\Delta x) = 0 \]
        
        \item Следовательно:
        \[ \lim_{\Delta x \to 0} \Delta f = 0 \]
        
        \item Равенство $\lim_{\Delta x \to 0} \Delta f = 0$ (или $\lim_{x \to x_0} f(x) = f(x_0)$) и является определением непрерывности функции в точке.
    \end{enumerate}
    \textbf{Вывод:} Функция непрерывна.
\end{tcolorbox}

\subsubsection*{Геометрический смысл производной и дифференциала}

\begin{center}
\begin{tikzpicture}[scale=1]
    % Оси
    \draw[->] (-0.5,0) -- (5,0) node[right] {$x$};
    \draw[->] (0,-0.5) -- (0,4) node[above] {$y$};
    
    % График функции
    \draw[thick] (0.5, 0.5) .. controls (1.5, 1) and (2.5, 0.5) .. (4.5, 3.5) node[right] {$y=f(x)$};
    
    % Точки
    \coordinate (M0) at (1.5, 0.85); % Точка касания x0
    \coordinate (M) at (3.5, 2.3);   % Точка секущей
    
    % Касательная
    \draw[red, thick] (0, -0.45) -- (4.5, 2.65) node[right] {Касательная $y_{kac}$};
    
    % Секущая
    \draw[blue, dashed] (0.5, 0.1) -- (4.5, 2.9) node[right] {Секущая};
    
    % Проекции
    \draw[dashed] (M0) -- (1.5, 0) node[below] {$x_0$};
    \draw[dashed] (M) -- (3.5, 0) node[below] {$x_0+\Delta x$};
    
    % Приращения
    \draw (1.5, 0.85) -- (3.5, 0.85);
    \draw (3.5, 0.85) -- (3.5, 2.3);
    
    % Треугольник дифференциала
    \draw[<->] (1.5, 0.5) -- (3.5, 0.5) node[midway, below] {$\Delta x$};
    
\end{tikzpicture}
\end{center}

\begin{tcolorbox}[title=Определение касательной]
    Пусть $k(\Delta x) \to k_0$ при $\Delta x \to 0$.
    Прямая $y = f(x_0) + k_0(x-x_0)$ называется \textbf{касательной} к графику функции $y=f(x)$ в точке $(x_0, f(x_0))$.
\end{tcolorbox}

\begin{tcolorbox}[title=Определение]
    \textbf{Касательная} — это предельное положение секущей, проходящей через точку касания $M_0$ и соседнюю точку $M$, когда $M \to M_0$.
\end{tcolorbox}

\textbf{Уравнение секущей:}
\[ \frac{y - f(x_0)}{x - x_0} = \frac{f(x_0+\Delta x) - f(x_0)}{\Delta x} \]

\begin{tcolorbox}[title=Геометрический смысл производной]
    Если функция $f(x)$ дифференцируема в точке $x_0$, то:
    \begin{enumerate}
        \item В точке $(x_0, f(x_0))$ существует невертикальная касательная к графику.
        \item Производная $f'(x_0)$ равна угловому коэффициенту этой касательной:
        \[ f'(x_0) = \tg \alpha = k \]
        где $\alpha$ — угол наклона касательной к положительному направлению оси $Ox$.
    \end{enumerate}
    \textbf{Следствие:}
    \begin{itemize}
        \item Если $f(x)$ непрерывна в $x_0$, то существует вертикальная касательная $\iff f'(x_0) = \infty$.
    \end{itemize}
\end{tcolorbox}

\begin{tcolorbox}[title=Геометрический смысл дифференциала]
    \textbf{Дифференциал} $df$ (или $dy$) есть приращение ординаты \textbf{касательной} к графику функции в данной точке.
    
    \[ \Delta f = f(x_0+\Delta x) - f(x_0) \quad \text{(приращение функции)} \]
    \[ df = f'(x_0)\Delta x = \tg \alpha \cdot \Delta x \quad \text{(приращение касательной)} \]
    
    При малых $\Delta x$: $\Delta y \approx dy$.
\end{tcolorbox}

\textbf{Доказательство касания:}
Нужно доказать, что расстояние $\rho$ от точки графика до касательной есть $o(\Delta x)$.
\[ \rho = |M_0 M| \cdot \sin(\dots) \]
Более строго: разность ординат графика и касательной:
\[ \Delta y - dy = (f(x_0) + A\Delta x + o(\Delta x)) - (f(x_0) + A\Delta x) = o(\Delta x) \]
\[ \lim_{\Delta x \to 0} \frac{\Delta y - dy}{\Delta x} = 0 \]
Это означает, что график функции и касательная сближаются быстрее, чем $\Delta x$.


